\section{Introduction}
\label{sec:intro}

% ── 1. Opening Context / Hook ───────────────────────────────────────────────
Large-language-model (LLM) agents are rapidly becoming a primary
interface between users and external services. By reasoning over tool
schemas and iteratively invoking APIs, these agents can complete
multi-step tasks---booking travel, synthesising medical records,
orchestrating data pipelines---that were previously tractable only
through hand-crafted workflows. The practitioner community has
converged on a small number of \emph{controller archetypes}:
\ReAct-style reasoning loops~\cite{yao2023react},
Toolformer-like inline API calls~\cite{schick2023toolformer},
and production frameworks such as LangChain's
\texttt{AgentExecutor}, the OpenAI Agents SDK, and AutoGen.
These controllers are mature and powerful, yet they share a
fundamental blind spot: they were designed for \emph{task success},
not for \emph{resource governance}.

% ── 2. Problem Statement ────────────────────────────────────────────────────
\paragraph{The gap.}
Real deployments impose hard constraints on agent behaviour: a
corporate policy may cap LLM API spend at $\$B$ per request; a
latency SLO for a customer-facing task may require a response within
$D$ seconds. Current controllers enforce these constraints only via
\emph{soft prompt instructions} of the form ``stay under budget'';
they have no mechanisms for formal guarantees. We identify three
concrete failure modes that arise even from careful prompt-based
governance:
\begin{enumerate}
  \item \textbf{Budget leakage.} Per-call costs follow heavy-tailed
    distributions; a single unexpected call can exhaust the budget
    before the agent checks it.
  \item \textbf{Deadline opacity.} Tool latency is non-stationary
    and right-skewed; prompt-based caps cannot account for
    tail-latency accumulation across sequential calls.
  \item \textbf{Discovery--cost coupling.} Large registry pools
    (e.g., MCP-style~\cite{mcp2024}) inflate both retrieval cost
    and mis-selection probability, feeding back into total spend.
\end{enumerate}

% ── 3. Motivation / Why It Matters ─────────────────────────────────────────
\paragraph{Why this matters.}
Budget overruns in agent deployments are not hypothetical. In
production LLM pipelines, a single runaway agent session can
consume orders of magnitude more tokens than budgeted, directly
translating into financial cost and SLO penalties. As enterprises
adopt agentic AI at scale, the absence of formal cost-safety
guarantees constitutes a critical governance gap that soft prompting
cannot close. This paper provides the first formal treatment of this
gap as a computational optimisation problem.

% ── 4. Related Work (Brief) ─────────────────────────────────────────────────
\paragraph{Prior approaches and their limitations.}
Service-oriented computing has a rich tradition of QoS-aware
composition~\cite{zeng2004qos,alrifai2009qos}, formulating
deployment as constrained optimisation over cost, latency, and
reliability. These methods assume deterministic or stationary
distributions and offline planning---assumptions that break down
under the stochastic, non-stationary cost of LLM-mediated tool
calls. In the online learning literature, \emph{bandits with
knapsacks} (\BwK)~\cite{bwk2018} provides the closest formal
framework, but it lacks (i)~schema-type filtering to prune
infeasible tools before exploration, and (ii)~explicit
chance-constraint semantics needed to bound violation probabilities.
Constraint-aware extensions---primal-dual Lagrangian
bandits~\cite{mahdavi2012tradeoff,agrawal2014bwk} and
CVaR-based formulations~\cite{tamkin2020cvar}---provide
theoretical tools but have not been instantiated for the
agent-orchestration setting. Benchmarks such as
ToolBench~\cite{qin2023toolllm}, $\tau$-bench~\cite{yao2024taubench},
and the Berkeley Function-Calling
Leaderboard~\cite{bfcl2024} provide rich task traces but do not
expose cost/latency ground truth in a form directly usable for
chance-constrained optimisation.

% ── 5. Your Approach / Key Idea ─────────────────────────────────────────────
\paragraph{Our approach.}
We introduce \CBSLAO~(\emph{Cost- and SLA-Bounded Agent
Orchestration}), a formal model in which an agent adaptively
selects tool calls to maximise expected utility subject to
\emph{chance constraints} on total cost and latency. The core
insight is that two complementary mechanisms---\emph{offline
schema-type pruning} to eliminate structurally infeasible tools
before any call is issued, and \emph{budget-aware UCB selection}
with a sub-Gaussian chance-constraint correction---can be composed
into a single policy, \CBUC, that provides both a formal regret
bound and strong empirical safety. Unlike Lagrangian or CVaR
relaxations, \CBUC~enforces constraints directly in the planning
step without requiring dual-variable tuning.

% ── 6. Contributions ────────────────────────────────────────────────────────
\paragraph{Contributions.}
\begin{itemize}
  \item[(C1)] \textbf{Problem formalisation.} A precise statement of
    \CBSLAO~with chance-constrained budget and deadline semantics,
    distinguishing adaptive, semi-adaptive, and non-adaptive policy
    classes (\S\ref{sec:formulation}).
  \item[(C2)] \textbf{Hardness.} NP-hardness of the decision version
    of \CBSLAO~by polynomial reduction from 0/1 Knapsack, together
    with inapproximability remarks. The reduction holds even for
    degenerate (point-mass) cost distributions
    (\S\ref{sec:hardness}).
  \item[(C3)] \textbf{Algorithm and regret bound.} \CBUC, an online
    policy with an $\tOh(\sqrt{KT})$ regret bound against the best
    non-adaptive feasible policy for $K$ tools over $T$ rounds
    (\S\ref{sec:algorithm}).
  \item[(C4)] \textbf{Empirical evaluation.} A reproducible simulator
    and controlled extended sweep ($3{,}960$ replicate rows) showing
    \CBUC~cuts budget-overrun from $36.9\%$ (\ReAct-capped) and
    $28.9\%$ (\BwK-vanilla) to $1.1\%$ and eliminates deadline misses,
    while making explicit the utility cost relative to Lagrangian and
    pre-check baselines (\S\ref{sec:evaluation},
    Appendix~\ref{app:ablations}).
\end{itemize}

% ── 7. Paper Organisation ───────────────────────────────────────────────────
\paragraph{Paper organisation.}
Section~\ref{sec:related} reviews related work in detail.
Section~\ref{sec:formulation} defines the \CBSLAO~problem formally.
Section~\ref{sec:hardness} presents the NP-hardness proof.
Section~\ref{sec:algorithm} describes \CBUC~and its regret analysis.
Section~\ref{sec:evaluation} reports experiments.
Sections~\ref{sec:threats}--\ref{sec:conclusion} discuss threats,
limitations, and conclusions.
Appendix~\ref{app:ablations} contains the three ablation studies.

\paragraph{Non-goals.}
We do not propose a new LLM, a new agent framework, or a production
runtime. \CBSLAO~is a \emph{governance-layer} contribution: an agent
framework can adopt \CBUC~as its sequencing policy without structural
changes to prompts, models, or tools.
